\begin{lemma}[Scheide symmetric defective internal edges (11)--(12)]
For each defective color $\delta$, there is a $\delta$-colored edge
identity internal to $U_{\gamma(\delta)}$, and there is a
$\gamma(\delta)$-colored edge identity internal to $U_\delta$.
\end{lemma}
\begin{proof}
Fix a defective color $\delta$.  Put
\[
 X=V(T),\qquad r=\gammaIndex(\delta),\qquad
 \gamma=\gamma(r),\qquad
 U_\gamma=\{u^1_\gamma,u^2_\gamma\},\qquad Z=X\setminus U_\gamma,
\]
and write $U_\delta=\{u^1_\delta,u^2_\delta\}$ for the replacement
pair.  All the vertices and edge identities in this notation are the
ordered data of \noderef{ScheideDefectiveSetup} and
\noderef{ScheideReplacementData}; in particular, the three distinct
$\delta$-edges $f^0_\delta,f^1_\delta,f^2_\delta$ have their tree ends in
$Z$, their other ends outside $X$, and are exactly the $\delta$-colored
departures having these two properties.  For $f^0_\delta$, the assertion
that its tree end lies in $Z$ follows as follows.  The setup says that
$v^0_\delta$ is a Scheide exit, and
\noderef{ScheideStructuralExitsInBase} puts every such exit in the base
$Y$.  The exact ordered layout in the underlying
\noderef{ScheideReducedConfiguration} and
\noderef{ScheideStructuralConfiguration}, together with the distinctness
in \noderef{TreeSequence}, gives $Y\subseteq Z$.

We first choose the free color used for (11), including the count.  The
inherited structural configuration says that $(T,e_0,\varphi)$ is
globally maximum.  Hence \noderef{ScheideGlobalMaximumOrder} says that
$X$ is elementary and closed.  Unpacking
\noderef{GlobalMaximumTriple} and \noderef{KTriple} also says that
$\varphi$ is proper in the sense of
\noderef{IsProperEdgeColoring}.  The structural layout gives
$|X|=p+2s$ and $|Z|=p+2s-2$.  The first edge of
\noderef{TreeSequence} is $e_0$ and its two distinct ends are the first
two base vertices, so both ends of $e_0$ belong to $Y\subseteq Z$.

Apply \noderef{ScheideGlobalMaximumMissingCounts} at every vertex of
$Z$.  Each vertex contributes at least one missing color, and each of the
two ends of $e_0$ contributes one further missing color.  Since $X$ is
elementary, the definition \noderef{ColoredVertexSetElementary} makes
all these missing colors distinct between different vertices.  Therefore
the union of the missing-color sets on $Z$ has cardinality at least
\[
 |Z|+2=p+2s. \tag{13}
\]
On the other hand, at most $p-2+s$ colors belong to
\noderef{TreeUsedColors}$(\varphi,T)$.  There are $p-2$ colored base
edges after $e_0$, and the two edges of each of the $s$ blocks have their
single common block color.  Every one of these used colors is missing at
a vertex of $Y$: this is explicit for a block color in
\noderef{ScheideStructuralConfiguration}, while for a colored base edge
it is the earlier-vertex witness in \noderef{TashkinovTree}, which is
part of the globally maximum triple and is still a base vertex at that
position.  Thus every used color is among the colors counted in (13).
Since
\[
 p-2+s<p+2s,
\]
there is a color $\beta$ missing at some vertex $v\in Z$ but unused on
the colored tree edges.  By \noderef{TreeFreeColors}, $\beta$ is free
for $(T,e_0,\varphi)$.  Elementariness makes $v$ unique.  Moreover
$\beta\ne\gamma$, because $\gamma$ colors the two edges of the
$r$-th block and hence is used on $T$.

Now instantiate \noderef{ScheideDefectiveChainPathExists} with the
globally maximum triple $(T,e_0,\varphi)$, defective color $\delta$,
free color $\beta$, and the vertex $v$ missing $\beta$.  It supplies an
outside vertex $z$ and a simple $(\delta,\beta)$ incidence path $P$ from
$v$ to $z$.  By \noderef{ScheideDefectiveChainPath}, the support of $P$
is the entire two-color component.  The conclusion of
\noderef{ScheideDefectiveChainBoundaryEdges} therefore puts every
$\delta$-colored boundary edge identity of $X$ on $P$.  In particular,
the three boundary identities $f^0_\delta,f^1_\delta,f^2_\delta$ occur
on $P$.

Apply \noderef{ScheideReplacementInternalMissing} to $\beta\ne\gamma$
and the fact that $\beta$ is missing in $Z$.  We obtain a
$\beta$-colored edge identity $h_\delta$ with ends
$u^1_\delta,u^2_\delta$.  The path order around these five edge and
vertex nodes is forced, not merely possible.  Indeed,
$f^i_\delta$ and $h_\delta$ are respectively the known incident
$\delta$- and $\beta$-edges at $u^i_\delta$.  Two distinct edges of one
color incident at one vertex would be adjacent by
\noderef{AdjacentEdges} and would contradict properness.  Hence those
are the unique incidences of the two path colors at that vertex.  This is
also the vertex-node case of
\noderef{TwoColorIncidenceDegreeAtMostTwo} in the incidence graph
\noderef{TwoColorIncidenceGraph}.  Each edge-node has only its two
endpoint incidences, and the endpoints of $h_\delta$ are distinct by
the loopless law in \noderef{LooplessMultigraph}.  Because $P$ is a
simple path whose support is the full component, it must consequently
contain, in one of the two orientations, the consecutive segment
\[
 f^1_\delta, u^1_\delta, h_\delta, u^2_\delta, f^2_\delta. \tag{14}
\]
Thus $f^1_\delta$ and $f^2_\delta$ are exactly one returning excursion
of $P$ from $Z$ through $U_\delta$ and back to $Z$.

Since $P$ begins in $X$ and ends outside $X$, it has a last vertex-node
$w$ belonging to $X$ in the order starting at $v$.  The defective color,
the free color missing at $v$, the full path $P$, and this last occurrence
are precisely the witnesses in \noderef{ScheideExitVertex}; hence $w$
belongs to the finite exit set of \noderef{ScheideExitVertices}.
By \noderef{ScheideStructuralExitsInBase}, $w\in Y\subseteq Z$.  Let
$q$ be the edge-node immediately after $w$ on $P$.  It is a boundary
edge of $X$ by \noderef{BoundaryEdges}.  It cannot have color $\beta$:
$\beta$ is missing in $X$,
and closedness, as defined in \noderef{ColoredVertexSetClosed}, forbids a
boundary edge with any color missing in $X$.  Thus $q$ has color
$\delta$.  Its tree end is in $Z$ and its other end is outside $X$, so
the exact-three clause of \noderef{ScheideReplacementData} gives
\[
 q\in\{f^0_\delta,f^1_\delta,f^2_\delta\}.
\]
Neither $f^1_\delta$ nor $f^2_\delta$ can be this final departure,
because (14) forces either one, when used as a departure, to return to a
later vertex of $Z\subseteq X$ through the other.  Therefore
$q=f^0_\delta$.  The endpoint record in
\noderef{ScheideDefectiveSetup} then identifies $w=v^0_\delta$.  This
establishes the claimed path order before using it below.

We next prove that $P$ avoids $U_\gamma$.  The lemma
\noderef{ScheideStructuralInternalPair}, applied to $\beta\ne\gamma$ and
the missing occurrence of $\beta$ in $Z$, supplies a $\beta$-edge
internal to $U_\gamma$.  Properness makes it the unique incident
$\beta$-edge at either vertex of this pair, so no $\beta$-edge joins
$U_\gamma$ to $Z$.  The assertion
\noderef{ScheideDefectiveNoCross} says independently that no
$\delta$-edge joins $U_\gamma$ to $Z$.

Suppose that $P$ nevertheless first reaches a vertex of $U_\gamma$.
There is no direct two-color incidence from $Z$ to $U_\gamma$, by the
preceding paragraph, so before that first occurrence the path must leave
$X$ from a vertex of $Z$.  Consider the last such departure before the
first occurrence of $U_\gamma$.  Closedness again says that its color is
$\delta$, and the exact-three clause says that it is one of
$f^0_\delta,f^1_\delta,f^2_\delta$.  If it is $f^1_\delta$ or
$f^2_\delta$, the forced segment (14) returns to $Z$ before the path can
go anywhere else, contradicting the choice of this departure.  If it is
$f^0_\delta$, it occurs after the last tree vertex just identified, so
the path cannot subsequently reach $U_\gamma\subseteq X$.  This exhausts
the three possibilities and proves
\[
 V(P)\cap U_\gamma=\varnothing. \tag{15}
\]

Every $\delta$-colored boundary identity of $X$ lies on $P$ by
\noderef{ScheideDefectiveChainBoundaryEdges}.  Hence (15) says that none
of them is incident with $U_\gamma$.  Also, because $\delta$ is defective,
\noderef{DefectiveColors} and \noderef{BoundaryColorEdges} give at least
one $\delta$-colored boundary identity of $X$.  Since $X$ is closed,
$\delta$ cannot be missing at any vertex of $X$; by the complementary
definitions \noderef{MissingColors} and \noderef{PresentColors}, it is
present at both vertices of $U_\gamma$.  Choose the incident
$\delta$-edge at either one.  Its other end cannot be in $Z$ by
\noderef{ScheideDefectiveNoCross}, and cannot be outside $X$ by (15) and
the boundary-path equality.  The loopless law excludes the same endpoint,
so the other end is the other vertex of $U_\gamma$.  Thus this individual
edge identity is internal to $U_\gamma$ and has color $\delta$, proving
(11).

For (12), set
\[
 \varphi_\delta=
 \operatorname{swap}_{\gamma,\delta}(\varphi;v^0_\delta),
 \qquad T_\delta=\text{the replacement tree},\qquad
 X_\delta=V(T_\delta)=Z\cup U_\delta.
\]
The coloring and the last equality are the exact data of
\noderef{ScheideDefectiveSwitchEffect} and
\noderef{ScheideReplacementData}.  The theorem
\noderef{ScheideReplacementMaximum} says that $\varphi_\delta$ is proper,
leaves exactly $e_0$ uncolored, and that
$(T_\delta,e_0,\varphi_\delta)$ is globally maximum.  Applying
\noderef{ScheideGlobalMaximumOrder} to this new triple shows explicitly
that $X_\delta$ is elementary and closed under $\varphi_\delta$.

We record exactly how $\gamma$ becomes defective.  In the original
coloring, $\gamma$ is missing at $v^0_\delta$, whereas the incident edge
$f^0_\delta$ has color $\delta$; by
\noderef{Incident}, \noderef{PresentColors}, and
\noderef{MissingColors}, this also proves $\delta\ne\gamma$.
The switch-effect theorem colors $f^0_\delta$ with $\gamma$ and fixes
every other edge identity incident with $X$.  Let
$g^1_\gamma,g^2_\gamma$ be the two edges of the deleted $\gamma$-block.
The structural configuration colors both with $\gamma$, and they remain
so under $\varphi_\delta$.  Their base ends lie in $Y\subseteq Z$ and
their block ends lie in $U_\gamma$, which is outside $X_\delta$.
Moreover $f^0_\delta$ also leaves $X_\delta$: its outside end cannot equal
$u^i_\delta$, since otherwise the distinct original $\delta$-edges
$f^0_\delta$ and $f^i_\delta$ would violate properness at that vertex.
Thus the three pairwise distinct identities
\[
 f^0_\delta,\quad g^1_\gamma,\quad g^2_\gamma \tag{16}
\]
are $\gamma$-colored boundary edges of $X_\delta$ with their tree ends
in $Z$.

They are exactly the $\gamma$-departures from $Z$ under
$\varphi_\delta$.  Indeed, let $a$ be such a departure.  If
$a\ne f^0_\delta$, the switch-effect theorem gives
$\varphi_\delta(a)=\varphi(a)=\gamma$.  Original closedness of $X$ and
the fact that $\gamma$ is missing at $v^0_\delta$ exclude an original
$\gamma$-edge from $Z$ to the exterior of $X$.  Hence the other end of
$a$ lies in $X\setminus X_\delta=U_\gamma$.  At either vertex of
$U_\gamma$, the corresponding one of $g^1_\gamma,g^2_\gamma$ is already
an incident $\gamma$-edge, and original properness makes it unique.
Therefore $a$ is that named block edge.  This proves the exactness claim.
In particular, the three identities in (16) show, via
\noderef{DefectiveColors}, that $\gamma$ is defective for
$X_\delta$ under $\varphi_\delta$.

We now perform the missing-color count for the replacement triple rather
than appealing to symmetry.  The two ends of $e_0$ still lie in
$Y\subseteq Z\subseteq X_\delta$.  Applying
\noderef{ScheideGlobalMaximumMissingCounts} to the replacement triple and
using its elementariness gives at least
$|Z|+2=p+2s$ distinct missing colors on $Z$, exactly as in (13).  At most
$p-2+s$ colors are used on $T_\delta$: there are at most $p-2$ retained
colored base-edge colors, $s-1$ retained block colors, and the two appended
edges share the single color $\delta$.  The last assertion follows from
\noderef{ScheideReplacementPairFixed}, since those appended edges are
incident with $U_\delta$ and originally have color $\delta$.

Every color used on $T_\delta$ has a missing witness in $Z$.  A retained
base or block color has its old witness in $Y$ as in the first count.  The
only missing set changed at an old tree vertex is that at
$v^0_\delta$, where \noderef{ScheideDefectiveSwitchEffect} replaces the
missing color $\gamma$ by $\delta$.  A retained base color differs from
every block color by \noderef{ScheideStructuralConfiguration}, and a
retained block color differs from $\gamma$ by injectivity of the indexed
block colors, so each retained witness survives.  The new used color
$\delta$ is missing at $v^0_\delta\in Y\subseteq Z$ by the same exact
switch formula.  Consequently the same strict cardinality inequality
chooses a color $\beta_\delta$ which is missing at a vertex $v_\delta\in Z$
and is unused on $T_\delta$.  By \noderef{TreeFreeColors}, it is free for
the replacement triple.  Since $\delta$ is used on the appended edges,
$\beta_\delta\ne\delta$.  Since $\gamma$ is defective and
$X_\delta$ is closed, $\gamma$ is present at every vertex of
$X_\delta$; hence $\beta_\delta\ne\gamma$.  Finally the switch-effect
formula, and these two inequalities, show that
$\beta_\delta$ is also missing under the original coloring at the same
vertex $v_\delta\in Z$.

Instantiate \noderef{ScheideDefectiveChainPathExists} for the globally
maximum replacement triple with defective color $\gamma$, free color
$\beta_\delta$, and missing vertex $v_\delta$.  We obtain a simple full
$(\gamma,\beta_\delta)$ incidence path $P_\delta$ from $v_\delta$ to an
outside vertex.  By \noderef{ScheideDefectiveChainPath} its support is
the full component, and by
\noderef{ScheideDefectiveChainBoundaryEdges} it contains every
$\gamma$-colored boundary identity of $X_\delta$, in particular all three
identities in (16).

The original missing occurrence of $\beta_\delta$ in $Z$ and
$\beta_\delta\ne\gamma$ allow us to apply
\noderef{ScheideStructuralInternalPair}.  It gives an original
$\beta_\delta$-edge $h_\gamma$ internal to $U_\gamma$.
The switch-effect theorem fixes this edge, because it is incident with
$X$ and is not $f^0_\delta$, so it is still colored $\beta_\delta$ under
$\varphi_\delta$.  Properness of $\varphi_\delta$, the degree bound
\noderef{TwoColorIncidenceDegreeAtMostTwo}, and the full-component
property now force the consecutive segment
\[
 g^1_\gamma, u^1_\gamma, h_\gamma,
 u^2_\gamma, g^2_\gamma \tag{17}
\]
in one of its two orientations on $P_\delta$, by exactly the incidence
degree argument used for (14).  Thus the two deleted block edges form one
returning excursion through $U_\gamma$.

We next locate the last replacement-tree vertex on $P_\delta$.  First,
every vertex $x\in U_\delta$ misses at least one color by
\noderef{ScheideGlobalMaximumMissingCounts} for the replacement triple.
Every used color of $T_\delta$ has a missing witness in $Z$, as just
proved.  Since $X_\delta$ is elementary and $x\notin Z$, none of those
used colors can be missing at $x$.  Therefore any color missing at $x$
is unused and hence is free by \noderef{TreeFreeColors}.  The theorem
\noderef{ScheideExitNoFree} says that a Scheide exit of a globally maximum
triple misses no free color.  It follows that neither vertex of
$U_\delta$ is an exit of the replacement tree.

Choose the last vertex-node $w_\delta$ of $X_\delta$ on $P_\delta$.
The definition \noderef{ScheideExitVertex} makes it a replacement-tree
exit, so the preceding paragraph and
$X_\delta=Z\cup U_\delta$ put it in $Z$.  Its next edge is a boundary
edge and cannot have color $\beta_\delta$, because $X_\delta$ is closed
and that color is missing on it.  It therefore has color $\gamma$.
The exact list (16) shows that this final departure is one of the three
named edges.  It is not $g^1_\gamma$ or $g^2_\gamma$, because (17) would
then return to a later vertex of $Z$.  It is consequently
$f^0_\delta$.  This is the replacement-coloring last-exit order needed
below.

There is also a $\beta_\delta$-edge $h'_\delta$ internal to $U_\delta$:
apply \noderef{ScheideReplacementInternalMissing} to the original
missing occurrence of $\beta_\delta$ in $Z$, and then use
\noderef{ScheideReplacementPairFixed} to retain its color under
$\varphi_\delta$.  Properness makes $h'_\delta$ the unique incident
$\beta_\delta$-edge at either vertex of $U_\delta$.  Furthermore no
$\gamma$-edge under $\varphi_\delta$ joins $U_\delta$ to $Z$.  If one
did, \noderef{ScheideReplacementPairFixed} would give the same
$\gamma$-colored identity under $\varphi$; since $U_\delta$ is outside
$X$ and $Z\subseteq X$, this would be an original $\gamma$-boundary edge
of $X$, contrary to original closedness and the fact that $\gamma$ is
missing at $v^0_\delta$.

Suppose $P_\delta$ reached $U_\delta$.  The internal edge
$h'_\delta$ and properness exclude a $\beta_\delta$ crossing from $Z$,
and the preceding paragraph excludes a $\gamma$ crossing.  Hence, before
its first occurrence in $U_\delta$, the path must leave $X_\delta$ from
$Z$.  Take the last such departure.  It has color $\gamma$ and hence is
one of (16).  If it is a deleted block edge, the forced excursion (17)
returns to $Z$ before proceeding; if it is $f^0_\delta$, it occurs after
the last replacement-tree vertex.  Both alternatives contradict the
choice of the departure.  Thus
\[
 V(P_\delta)\cap U_\delta=\varnothing. \tag{18}
\]

All $\gamma$-colored boundary edges of $X_\delta$ lie on $P_\delta$, so
(18) excludes every such boundary edge incident with $U_\delta$.  Since
$\gamma$ is defective and $X_\delta$ is closed, the complementary
definitions \noderef{MissingColors} and \noderef{PresentColors} show
that $\gamma$ is present at both vertices of $U_\delta$.  At either
vertex choose its unique incident $\gamma$-edge.  Its other end is not in
$Z$ by the no-cross assertion above and not outside $X_\delta$ by (18)
and the boundary-path equality.  Looplessness leaves only the other
vertex of $U_\delta$.  We have therefore found an individual edge
identity internal to $U_\delta$ and colored $\gamma$ under
$\varphi_\delta$.  It is incident with $U_\delta$, so
\noderef{ScheideReplacementPairFixed} transfers the same color to the
original coloring $\varphi$.  This proves (12), and hence both asserted
internal-edge conclusions.
\end{proof}
